\section{Weak forms on the solid reference configuration}
\label{sec:reference_solid_weak_forms}

Let \(\test{\mbf u}\), \(\test q_f^\alpha\), and \(\test q_s\) denote
finite-element test functions on \(\refdomain\).  Boundary subsets with
prescribed traction or flux are denoted by \(\Gamma_{0t}\) and
\(\Gamma_{0W}\).  The reference outward unit normal is \(\mbf N\).

\subsection{Fluid component balance}
\label{sec:fluid_component_weak_form}

For component \(\alpha\) in fluid phase \(f\), define the solid-reference
storage
\[
	M_f^\alpha = J\phi_f\bar\rho_f\eta_f^\alpha .
\]
The strong form on the solid trajectory is
\[
	\left.\frac{\partial M_f^\alpha}{\partial t}\right|_{\mbf X}
	+\Div\left(\eta_f^\alpha\mbf W_f\right)
	=
	J\sum_m \nu_{f(m)}^\alpha \dot r_{(m)} .
\]
The weak form uses the same reference measure; no current-volume source term is
inserted without the explicit pull-back factor \(J\).
The Galerkin residual is
\begin{align}
	R_f^\alpha
	=&
	\int_{\refdomain}
	\test q_f^\alpha
	\left[
		\left.\frac{\partial M_f^\alpha}{\partial t}\right|_{\mbf X}
		-
		J\sum_m \nu_{f(m)}^\alpha \dot r_{(m)}
	\right]\dd V
	\nonumber \\
	&-
	\int_{\refdomain}
	\Grad \test q_f^\alpha \cdot
	\left(\eta_f^\alpha\mbf W_f\right)\dd V
	+
	\int_{\Gamma_{0W}}
	\test q_f^\alpha
	\eta_f^\alpha\bar W_f
	\dd A .
	\label{eq:fluid_component_residual}
\end{align}
Here \(\bar W_f=\mbf W_f\cdot\mbf N\) is the prescribed outward reference mass
flux of the phase before multiplication by \(\eta_f^\alpha\).  The sign
convention makes outward flux positive.

\subsection{Solid component or solid phase balance}
\label{sec:solid_component_weak_form}

For a one-component solid phase \(s\), the storage variable is
\[
	M_s = J\phi_s\bar\rho_s .
\]
Because each solid phase is attached to the skeleton, there is no relative-flux
term:
\begin{equation}
	R_s =
	\int_{\refdomain}
	\test q_s
	\left[
		\left.\frac{\partial M_s}{\partial t}\right|_{\mbf X}
		-
		J\sum_m \nu_{s(m)}^s \dot r_{(m)}
	\right]\dd V .
	\label{eq:solid_component_residual}
\end{equation}
If several solid components share a solid phase, this residual must be promoted
to a component balance with the appropriate mass fractions and conversion
sources.

\subsection{Fluid phase kinematic residuals}
\label{sec:phase_kinematic_weak_forms}

No mobile-phase deformation-gradient or Jacobian residual is part of the first
solid-reference finite-element system.  Phase-attached scalar derivatives are
evaluated with \eqref{eq:fluid_material_derivative_on_solid_frame}, using the
same \(\mbf W_f\) that enters the component balance.  If a later closure
requires \(J_f\) or the full tensor \(\mbf F_f\), their evolution should be
introduced as auxiliary/history kinematics and tied back to the phase map at
that point.

\subsection{Reference flux closure}
\label{sec:reference_flux_closure}

The reference relative mass flux entering
\eqref{eq:fluid_component_residual} has the form
\begin{align}
	\mbf W_f
	= & \;
	\rho_f\widetilde{\mbf K}_f
	\left[
		\rho_f\mbf F^T\left(\mbf g-\mbf a_f\right)
		-\rho_f
		\left(
			\Grad\psi_f
			+
			\mathfrak{s}_f\Grad\theta_{\mc F}
		\right)
		-\phi_f\Grad
		\left(
			\bar p_E+\gamma_f
		\right)
		\right.
		\notag \\
	& \phantom{\;\rho_f\widetilde{\mbf K}_f}
		\left.
		+
		\mbf F^T
		\left[
			\grad\cdot
			\left(
				\phi_f\mbf E\otimes\mbf d_f
				\right)
		\right]
		\right.
		\notag \\
	& \phantom{\;\rho_f\widetilde{\mbf K}_f}
		\left.
		+
		\sum_\alpha\dot c_f^\alpha
		\left(
			\Grad \tau
			-\mbf F^T\mbf v_s
		\right)
	\right],
	\label{eq:reference_flux_closure}
\end{align}
where \(\rho_f=\phi_f\bar\rho_f\) and
\begin{equation}
	\widetilde{\mbf K}_f
	=
	J\mbf F^{-1}
	\left(
			\phi_f^2\mu_f^{\mathrm{visc}}\bs\kappa_f^{-1}
		+\sum_\alpha\dot c_f^\alpha\mbf I
	\right)^{-1}
	\mbf F^{-T}.
	\label{eq:reference_modified_permeability}
\end{equation}
This is the Piola transform of the current generalized Darcy mass flux.  The
tensor \(\widetilde{\mbf K}_f\) is a pulled-back
conversion-corrected mobility, not a permeability: its current-configuration
inverse contains the net conversion phase-mass source
\(\sum_\alpha\dot c_f^\alpha\).  The inverse exists only when none of the
eigenvalues
	\(\phi_f^2\mu_f^{\mathrm{visc}}/\kappa_{f,i}+\sum_\alpha\dot c_f^\alpha\) vanishes.
Positive definiteness is the stronger positive-resistance condition
\(\sum_\alpha\dot c_f^\alpha>
	-\phi_f^2\mu_f^{\mathrm{visc}}/\kappa_{f,\max}\); it is sufficient, but not necessary, for
algebraic elimination.  These source terms and the \(\Grad\tau\) insertion
force are part of the conversion-momentum structure and must not be dropped
when reactions or phase transformations produce appreciable mass in a mobile
phase.  The additional
\(-\rho_f(\Grad\psi_f+\mathfrak{s}_f\Grad\theta_{\mc F})\) term is the
reversible bulk-transport contribution required by the theory paper's complete
fluid--skeleton interaction closure; it is not part of the dissipative drag
tensor.

\subsection{Transfer-potential and transfer-work residuals}
\label{sec:conversion_potential_weak_form}

Let \(\test t\) be the scalar test function for the single transfer potential
\(\tau\).  The Galerkin residual associated with
\eqref{eq:implementation_tau_solid_frame} is
\begin{align}
	R_\tau
	=&
	\int_{\refdomain}
	\test t
	\left[
		\left.\frac{\partial\tau}{\partial t}\right|_{\mbf X}
		-\frac{1}{2}\mbf v_s\cdot\mbf v_s
		-\hat\mu_{s_\star}^N
		+\psi_{s_\star}
		+\frac{\pi_{s_\star}^N}
		{\phi_{s_\star}\bar\rho_{s_\star}\eta_{s_\star}^N}
	\right]\dd V .
	\label{eq:tau_evolution_residual}
\end{align}
For every non-reference phase \(\xi\ne s_\star\), let
\(\test\ell_\xi\) test its algebraic reference-component transfer-work level.
Using \eqref{eq:fluid_material_derivative_on_solid_frame} for a mobile phase,
the recovery residual is
\begin{align}
	R_{L_f^N}
	=&
	\int_{\refdomain}
	\test\ell_f
	\left[
		L_f^N
		-\left.\frac{\partial\tau}{\partial t}\right|_{\mbf X}
		-\frac{\mbf W_f}{J\rho_f}\cdot\Grad\tau
		+\frac{1}{2}\mbf v_f\cdot\mbf v_f
		-\frac{\pi_f^N}
		{\phi_f\bar\rho_f\eta_f^N}
	\right]\dd V .
	\label{eq:fluid_transfer_work_recovery_residual}
\end{align}
For a non-reference solid phase attached to the skeleton, the corresponding
algebraic residual is
\begin{align}
	R_{L_s^N}
	=&
	\int_{\refdomain}
	\test\ell_s
	\left[
		L_s^N
		-\left.\frac{\partial\tau}{\partial t}\right|_{\mbf X}
		+\frac{1}{2}\mbf v_s\cdot\mbf v_s
		-\frac{\pi_s^N}
		{\phi_s\bar\rho_s\eta_s^N}
	\right]\dd V .
	\label{eq:solid_transfer_work_recovery_residual}
\end{align}
The traditional mechanism affinity
\eqref{eq:implementation_affinity_projection} is then an algebraic material
property or constraint residual.  If chemical or electrochemical equilibrium
is enforced directly, the residual is
\begin{equation}
	R_{\mc A_{(m)}}=
	\int_{\refdomain}
	\test a_{(m)}
	\left[
		\mc A_{(m)}
		+\sum_\xi\sum_\alpha
		\mu_\xi^\alpha
		\nu_{\xi (m)}^\alpha
	\right]\dd V ,
	\label{eq:affinity_projection_residual}
\end{equation}
with \(\test a_{(m)}\) the test function for an affinity multiplier or local
constraint variable.  A kinetic phase-change or pore-filling mechanism is
instead driven by the full conversion coefficient in
\eqref{eq:implementation_generalized_conversion_affinity_relation}; replacing
that coefficient by \(\mc A_{(m)}\) is valid only when the phasewise
stoichiometric mass sums vanish or the transfer-work offsets have equilibrated.

\subsection{Overall momentum balance}
\label{sec:momentum_weak_form}

Let \(\mbf P^{\prime\prime}\) be the skeleton effective first
Piola--Kirchhoff stress and \(B\bar p_E J\mbf F^{-T}\) the equivalent pore
pressure contribution.  The quasi-static implementation target is
\begin{align}
	R_u
	=&
	\int_{\refdomain}
	\Grad \test{\mbf u} :
	\left(
		\mbf P^{\prime\prime}
		-
		B\bar p_E J\mbf F^{-T}
	\right)\dd V
	-
	\int_{\refdomain}
	\test{\mbf u}\cdot J\rho\mbf g\dd V
	\nonumber \\
	&-
	\int_{\refdomain}
	\test{\mbf u}\cdot J
	\left[
			\sum_f\sum_\alpha
			\dot c_f^\alpha
			\left(
				\grad\tau-\mbf v_f
			\right)
			+
			\sum_s
			\dot c_s^s
			\left(
				\grad\tau-\mbf v_s
			\right)
	\right]\dd V
	\nonumber \\
	&-
	\int_{\Gamma_{0t}}
	\test{\mbf u}\cdot \bar{\mbf t}_0\dd A .
	\label{eq:momentum_residual}
\end{align}
For low-speed reservoir applications the inertial term from the full momentum
balance is omitted.  If retained, the term
\(\int_{\refdomain}\test{\mbf u}\cdot J\sum_\xi\rho_\xi\mbf a_\xi\dd V\)
is added to \eqref{eq:momentum_residual} with the same sign convention as the
chosen time integrator.  In a fully reference-form implementation, the
transfer-potential gradient in the conversion-momentum terms is evaluated as
\(\grad\tau=\mbf F^{-T}\Grad\tau\) for \(\tau\) stored on the solid mesh, while
\(\mbf v_f\) is reconstructed from
\eqref{eq:implementation_fluid_velocity_reconstruction}.

\subsection{Constraint residuals}
\label{sec:constraint_residuals}

Volume fractions and phase/component mass fractions satisfy
\[
	\sum_\xi \phi_\xi = 1, \qquad
	\sum_\alpha \eta_\xi^\alpha = 1 .
\]
The implementation may enforce these constraints by eliminating dependent
variables, by solving pressure or potential multipliers, or by adding explicit
constraint residuals.  The selected strategy must preserve the thermodynamic
conjugacy used to define \(\bar p_f\), \(\bar p_E\), \(B\), and
\(\mu_\xi^\alpha\).  This is a design constraint, not a numerical preference.
